\begin{abstract}
Federated learning is vulnerable to model-poisoning attacks because the server must aggregate client updates without inspecting client data. Existing robust aggregation rules often rely on update geometry, coordinate-wise robust statistics, or a single trusted reference direction, but these signals can confuse malicious behavior with benign client heterogeneity under non-i.i.d.\ data. We propose \textbf{CARAT}, a class-aware robust aggregation rule that uses labeled server-side examples to sample hidden, class-balanced tasks and score candidate updates by task-level loss reduction. CARAT combines these certificates with robust clipping, common-radius evaluation, coordinate-rank anomaly penalties, and capped-simplex weight optimization. We evaluate CARAT over five seeds on CIFAR-100/ResNet18 with $20\%$ malicious clients, three non-i.i.d.\ regimes, and five attacks. For $\alpha\in\{1,0.5\}$, CARAT has average rank $2.5$ under four conventional attacks and remains within $0.1$--$1.2$ percentage points of the strongest mean accuracy. Under Mimic, it ranks second among defenses with recorded results. At $\alpha=0.1$, however, CARAT becomes highly seed-sensitive under FangAttack and MinSum, exposing a clear robustness boundary. A component study further finds only $0.06$--$0.28$ percentage-point separation between the complete method and reduced variants. The evidence therefore supports conditional competitiveness, not universal superiority or large independent gains from every component.
\end{abstract}
